\section{1999: Sixteen Priests, One Protocol, and a Suspended Chapter}

\subsection{A warning about the evidence in this section and the next}

For the years 1999 and 2000 the documents that decide the question are
Roman acts addressed to the Fraternity: a letter of the Congregation for
Divine Worship of 3 July 1999 bearing protocol number 1411/99, a letter of
the Pontifical Commission \work{Ecclesia Dei} of 29 June 2000 numbered
748/2000, and a decree of the same Commission of 12 October 2000 numbered
823/2000. None of the three was published by the Holy See. None is in
\work{Acta Apostolicae Sedis}, on vatican.va, or in the Holy See Press
Office bulletin, and none was located in any official
publication.\endnote{Bounded negative, checked 2026-07-25. The 1999 letter
has an official printed locus that this account could not reach:
Congregatio de Cultu Divino et Disciplina Sacramentorum, \latin{Responsa
officialia} beginning \latin{Post liturgicam}, 3 July 1999, \work{Notitiae}
35 (1999) 307--309, as cited by Pietras, art.\ cit., 181 and 179 n.~23.
\work{Notitiae} is not distributed free online and the 1999 fascicle was
not obtained. The 2000 letter and decree have no printed locus known to
this account.}

What exists instead is a set of French, Italian and German texts of those
acts published between 1999 and 2001 on partisan websites---one an archive
of a French lay traditionalist association, one an Italian traditionalist
site---and, in 2026, cited from the same Italian site, with the German
original, by a peer-reviewed canon-law article whose author is a priest of
the Fraternity. That convergence is worth something: it establishes that
the texts circulating on those archives are the texts canonists work
from. It does not make them official, and this account never treats them
as official. Where a clause of one of these acts carries weight below, the
text is named, the language is named, and the host is named.

\subsection{The recourse of 29 June 1999}

The Fraternity's constitutions, approved \latin{ad experimentum} in 1988,
provided for a general chapter in 2000. In 1999 the superior general
obtained a dispensation to bring part of that chapter forward, so that
the definitive text of the constitutions could be prepared a year early;
elections were held in the spring to choose the capitulants.

The elections returned, on the account of the losing side, a body
composed entirely of priests opposed to any adaptation of the 1962 rite.
On 29 June 1999---the feast of Saints Peter and Paul---sixteen priests of
the Fraternity addressed a letter to Cardinal Angelo Felici, president of
the Pontifical Commission \work{Ecclesia Dei}, over the head of their own
superior general. Its first paragraph concedes the gravity of what they
were doing: ``In addressing ourselves directly to you without passing
through the authority of our Superior General, we are conscious of
performing a grave act, contrary to clerical habits.''\endnote{``Recours
de 16 prêtres de la FSSP,'' dated Rome, 29 June 1999, addressed to
Cardinal Angelo Felici, French text published in the archive of the
Association A.M.D.G.\ (amdg.asso.fr, ``Situation de la Fraternité Saint
Pierre 1999--2001''), read 2026-07-25. A partisan lay archive publishing
what it presents as the signatories' text; no official publication
exists, and this account uses it for the signatories' own positions and
requests only. The archive names the signatories with their apostolates.}

The letter's grievances are specific and worth reporting exactly, because
the Roman answer tracks them. It complains that seminarians refuse with
impunity to serve the Masses of visiting professors who celebrate the
\work{Novus Ordo Missae}; that congregational singing of the Our Father,
admitted a few years earlier, is now strictly forbidden; that a
liturgical directory adopted in 1995 with an appendix of adaptations for
France was never presented to seminarians; that concelebration with the
bishop ``in whatever circumstances'' is treated as an attack on the
Fraternity's unity and a grave delict, that members who concelebrated have
been sanctioned, and that before tonsure candidates must now undertake
not to concelebrate under pain of exclusion. It alleges a dysfunction in
the superior general's exercise of power, and it reports that when a
bishop offered territorial parishes subject to some indispensable local
liturgical adjustments, the offer was refused ``by fidelity to the
Constitutions.''

Its diagnosis is the sentence that would be quoted against the Fraternity
for years: that this liturgical hardening is not really a liturgical or
even a doctrinal position but ``rather psychological, even sociological,''
and that it repeats step by step the slide toward a spirit of separatism
that led the Society of Saint Pius X to refuse the Roman offers in 1988.
And it asks for three things: postponement of the chapter set for that
summer; the sending of a canonical visitor for the whole Fraternity; and
``the urgent appointment of an apostolic administrator to take in hand the
destiny of our Society.''

\subsection{Protocol 1411/99, 3 July 1999}

Four days later the Congregation for Divine Worship and the Discipline of
the Sacraments issued a document headed ``Official responses,'' protocol
1411/99, signed by Cardinal Jorge Medina Estévez, prefect, and Archbishop
Francesco Pio Tamburrino, secretary. It answers three questions put by an
unnamed priest of a community enjoying the faculty of the older rite. Its
premise is stated in its own preamble: because the use of the
pre-conciliar Missal is conceded only by indult, the liturgical right in
favour of the common Roman rite remains, and the Missal in force is the
one promulgated after the Council.

The three answers, in the French text available to this account, are
these:\endnote{``Protocole 1411,'' Congrégation pour le Culte Divin,
``Réponse officielle,'' 3 July 1999, French text in the A.M.D.G.\ archive
(amdg.asso.fr), read 2026-07-25. Unofficial French rendering; the Latin
original was not obtained. Its official printed locus is \work{Notitiae}
35 (1999) 307--309 (see the note above), which this account could not
reach; the substance summarized here is corroborated in outline by
Fr Bisig's own circular of 30 August 1999 and by the German text of the
letter of 29 June 2000, both discussed below. The paraphrase is this
account's and no project translation is offered.}

\begin{enumerate}
\item A priest of such a community \emph{may} freely use the Missal of
Paul VI when he celebrates for the good of the faithful, even
occasionally, in a parish where that Missal is used; indeed he
\emph{must} celebrate according to the post-conciliar Missal if the
celebration takes place in a community that uses the current Roman rite,
so that no astonishment or unease arises among the faithful.
\item Superiors of such communities, ``of whatever dignity,'' \emph{may
not} forbid the priests of their institutes the use of the post-conciliar
Missal when they celebrate for the good of the faithful in a community
where that Missal is used, because the indult is granted in the interest
of the faithful attached to the older rite and cannot be imposed on
communities that celebrate otherwise---over which, moreover, such
superiors have no authority.
\item A priest of such an institute \emph{may} concelebrate a Mass in the
current Roman rite; a superior or ordinary ``neither can nor should''
forbid him to concelebrate; on the contrary it is praiseworthy that he
concelebrate freely, above all at the Mass of Holy Thursday over which the
diocesan bishop presides, the sign of communion at the chrism Mass being
so strong that it should never be renounced save for a grave reason.
\end{enumerate}

The legal architecture of that answer deserves attention, because it is
the architecture that failed a year later. The responses treat the older
books as an \emph{indult}---a concession from the common law---and
therefore reason that the common law survives underneath, unimpaired, as a
right personal to each priest. On that premise the conclusion is
inescapable: what the law gives to every Latin priest, no superior below
the Roman Pontiff can take away.

\subsection{The Commission's measures of 13 July 1999}

Ten days later the Pontifical Commission acted on the sixteen priests'
first request and part of the third. According to the contemporary
account in the same archive, the Commission cancelled the general chapter
session convoked for August 1999; convoked instead an assembly of all
members incardinated in the Fraternity for the autumn; and restricted the
superior general to ordinary business, forbidding any change that was not
strictly necessary. That last measure had a practical bite the archive
notes: it protected the sixteen signatories from being
transferred.\endnote{``Situation à la Fraternité Saint Pierre 1999--2001''
and ``Données sur la crise au sein de la Fraternité Saint Pierre,''
A.M.D.G.\ archive (amdg.asso.fr), read 2026-07-25, a contemporary partisan
narrative. Its central claims about July 1999 are independently
corroborated by the superior general's own circular of 30 August 1999,
next cited, which reports the cancellation of the chapter session and the
substitute plenary meeting; the restriction on his powers is corroborated
indirectly by the Fraternity's own communiqué of 12 February 2000, which
states that ``the Superior general recovers the full faculties which the
law of the Church grants him.'' Neither the Commission's letter of 13 July
1999 nor its text was located.}

\subsection{The superior general's circular, 30 August 1999}

Father Bisig wrote to the members on 30 August 1999. The letter is
available to this account in a French text published, with evident
polemical intent, by the French district of the Society of Saint Pius X.
It is nevertheless the closest thing to a first-person Roman-facing
account of the crisis from the Fraternity's side, and its restraint is
itself evidence.

He names the two events: ``the Letter of 3 July 1999 (protocol 1411/99)
of the Congregation for Divine Worship, which presents itself as a
response to certain questions concerning the use of the \work{Novus Ordo
Missae} by priests attached to the traditional Roman liturgy,'' and ``the
cancellation by the Ecclesia Dei Commission of the session of the
Fraternity's general Chapter, which was to meet in August,'' replaced by
a plenary meeting of incardinated members in November. He states that a
small group of members lodged an official recourse with the Commission on
29 June ``without my knowledge.'' He states that he has lodged a recourse
of his own within the prescribed time, and gives his reason in a sentence
that states the whole constitutional question of this history: ``It is
important that the Fraternity rapidly recover the exercise of the
legislative and executive rights attached to its status as an Institute of
pontifical right.''\endnote{``Lettre de l'abbé Bisig aux membres de la
Fraternité Saint-Pierre---30 août 1999,'' French text published by the
District of France of the Society of Saint Pius X (laportelatine.org),
read 2026-07-25. A hostile party publishing an internal circular of the
institute it opposes; no original was located and no translator is named.
The quotations above render the published French closely and are given as
paraphrase. The letter is used for the superior general's own account of
the two Roman measures and for his stated legal objection.}

He also states his reasons for having refused the pressures to
concelebrate: ``fidelity to our proper mission in the Church, the concern
to avoid troubles and divisions among our faithful and our members, the
consolidation and identity of our apostolate, the unity of life and
discipline according to our Constitutions.'' And he offers a
characterization of the Fraternity's purpose that has some force: ``our
particular mission in the service of the hierarchy and of the faithful is
today more qualitative than quantitative''---a witness to the immemorial
liturgical tradition, ``sometimes to the detriment of the expansion of our
apostolate.''

\subsection{The joint supplication of 23 July, and Rome's public line}

On 23 July 1999 Father Bisig and Father Louis-Marie de Blignières, prior
general of the Fraternity of Saint Vincent Ferrer, jointly petitioned the
Holy See not to publish the responses of 3 July, on three grounds: that
they respect neither the proper character of the institutes nor the
jurisdiction of their superiors; that they open the way to a habitual
biritualism; and that they make the government of the institutes
impossible. The archive that publishes the supplication records that as
of 23 September 1999 it had received no answer.\endnote{Same A.M.D.G.\
archive, ``Données sur la crise,'' which summarizes the supplication under
its date and states the three grounds. The text of the supplication itself
was not located. The document does not say to which organ it was
addressed.}

Publicly, the Commission's line was that nothing extraordinary had
occurred. A text distributed by Monsignor Perl at a meeting of Una Voce
International in November 1999, and circulated thereafter, states that the
Commission ``has never had the intention of modifying'' the Fraternity's
statutes; that documents which should have remained strictly confidential
had been published, to the Fraternity's great harm; that several priests
had exercised ``their right to appeal to the Holy See---a right which every
faithful Catholic possesses''; and that those who call the resulting
measures an abuse of power ``do not understand the true juridical
situation: namely that the Commission exercises the full authority of the
Holy See over the aforementioned Fraternity.''\endnote{``Éclaircissements
fournis par Mgr Perl,'' French translation of an English text, A.M.D.G.\
archive (amdg.asso.fr), read 2026-07-25, which describes the text as read
by Monsignor Camille Perl at a meeting of Una Voce International in Rome
in November 1999 and reproduces it from a sheet distributed there. A
translation of a transcription on a partisan archive: the weakest evidence
used anywhere in this account, and it is used only for the one point that
is corroborated by an official act---the Commission's claim to exercise the
full authority of the Holy See over the Fraternity, which is exactly what
faculty 6 of \latin{Quia peculiare munus} confers.}

That last sentence is the pivot. Whatever else was disputed, both sides
agreed that the Pontifical Commission exercised the authority of the Holy
See over the Fraternity. The dispute was about what that authority could
reach.

\actrecord{Congregation for Divine Worship, ``Official responses,''
Prot.\ 1411/99, 3 July 1999 (French text on a partisan archive; official
locus \work{Notitiae} 35 [1999] 307--309, not reached).}{That in July 1999
the competent liturgical dicastery answered, on the premise that the 1962
books are held by indult, that a priest of an \work{Ecclesia Dei}
community may and in some circumstances must celebrate with the Missal of
Paul VI; that his superior may not forbid him to do so or to
concelebrate; and that concelebration at the chrism Mass is
praiseworthy.}{A responsum of a congregation, reached only in an
unofficial French rendering hosted by an interested party. It concerns
liturgical law and personal rights; it does not purport to interpret the
Fraternity's constitutions, name the Fraternity, or alter its erection.}
